\documentclass[../main.tex]{subfiles}
\graphicspath{{\subfix{../images/}}}
\begin{document}
	Analisi in Frequenza di Circuiti in Regime Sinusoidale.\\
	Un circuito che include componenti reattivi ha una risposta che cambia a seconda della frequenza in entrata.\\
	Le trasformate di condensatori e induttanze sono:
	\begin{gather}
		\bar z_L = j\omega L \\
		\bar z_C = \frac{1}{j\omega C}
	\end{gather}
	La velocità angolare $\omega$ viene espressa in radianti al secondo$\left[ rad/s \right]$, ed è uguale a $\omega = 2\pi f$, dove $f$ rappresenta la frequenza in cicli al secondo (Hertz $\left[ Hz \right]$).\\\\
	Un filtro è un circuito che lascia passare una certa gamma di frequenze, e si distingue in ideale o reale a seconda di quanto può essere preciso (un filtro ideale annulla completamente le altre frequenze, mentre un filtro reale le attenua).\\
	Quello che è importante capire nei filtri è che la frequenza rimane invariata, mentre vengono alterati ampiezza e fase del segnale in uscita.
	\subsection{Diagramma di Bode in Ampiezza e Fase}
	Un diagramma di Bode è una rappresentazione della risposta in frequenza di un sistema, come in fig.\ref{fig:im-26}. \\
	Esso indica come vengono modificati ampiezza e fase del segnale data la frequenza in entrata.
	\subsection{Filtri}
	\subsubsection{Filtro RC}
	Si tratta di un filtro formato da una resistenza e un condensatore in serie.\\
	La tensione in entrata è data da:
	\begin{equation}
		e(t)=\sqrt{2}\cos(\omega t)	
	\end{equation}
	\begin{center}
		\begin{figure}[h]
			\centering
			\includegraphics[scale=0.55]{example-image-a}
			\caption{Circuito filtro RC}
			\label{fig:im-25}
		\end{figure}
	\end{center}
	Costruendo l'equivalente di Steinmetz si può scrivere la tensione che si può prelevare ai capi del condensatore:
	\begin{gather}
		\bar V_C = \bar I \cdot \frac{1}{j\omega C} = \frac{E}{1+j\omega RC} \\
		|\bar V_C| = \frac{E}{\sqrt{1+\omega^2R^2C^2}}
	\end{gather}
	Il filtro risulta essere un \textbf{passa-basso} (considerando la tensione ai capi del condensatore) poiché se la frequenza tende a zero $\bar V_C\rightarrow E$, mentre se tende all'infinito $\bar V_C \rightarrow 0$. \\
	Le frequenze basse non vengono attenuate, mentre quelle alte si. Da notare che la curva di attenuazione procede in modo naturale e senza discontinuità (fig.\ref{fig:im-26}).
	\begin{center}
		\begin{figure}[h]
			\centering
			\includegraphics[scale=0.55]{example-image-a}
			\caption{Diagramma di Bode filtro RC passa-basso}
			\label{fig:im-26}
		\end{figure}
	\end{center}
	La frequenza di taglio, quindi il valore della velocità angolare in cui si dimezza la potenza (che è legata al quadrato della corrente) si ottiene:
	\begin{gather}
		\omega_C = \frac{1}{RC} 
	\end{gather}
	Dato che la resistenza non è un componente reattivo, la tensione ai capi di essa è in fase con il generatore, mentre per quello che riguarda la tensione ai capi del condensatore, essa viene sfasata con un angolo di 90º (vedi diagramma vettoriale di fig.\ref{fig:im-27}).
	\begin{center}
		\begin{figure}[h]
			\centering
			\includegraphics[scale=0.55]{example-image-a}
			\caption{Diagramma vettoriale}
			\label{fig:im-27}
		\end{figure}
	\end{center}
	Geometricamente la presenza dell'angolo retto implica che tra lo zero, $\bar V_R$ e $\bar E$ esiste una circonferenza, e se si varia il valore del condensatore $\bar V_R$ continua a rimanere su questa circonferenza. \\
	Quando ci si trova sul valore della pulsazione critica $\omega = \omega_C = \frac{1}{RC}$ la corrente diventa con parte reale e immaginaria uguali:
	\begin{equation}
		\bar I = \frac{E}{R+\frac{1}{j\frac{1}{RC}}C}= E \frac{R+jR}{R^2} = \frac{E}{R} + j\frac{E}{R}	
	\end{equation}
	Dato che le tensioni $\bar V_R$ e $\bar V_C$ sono uguali, l'angolo formato risulta $-\frac{\pi}{4}$.\\\\
	Facendo tendere la frequenza a zero, si nota che la tensione ai capi del condensatore tende a $\bar E$, mentre quella sulla resistenza diventa infinitesima, viceversa se $\omega \rightarrow +\infty$ (l'impedenza tende ad annullarsi).\\\\
	Se nello stesso circuito si prendesse la tensione in uscita sulla resistenza si avrebbe una situazione diversa:
	\begin{equation}
		\bar V_R = R\cdot \bar I = R\cdot \frac{\bar E}{R+\frac{1}{j\omega C}} = E \frac{j\omega RC}{j\omega RC + 1}	
	\end{equation}
	Il modulo della tensione è quindi:
	\begin{equation}
		|\bar V_R| = \frac{E\omega RC}{\sqrt{\omega ^2 R^2 C^2 + 1}}	
	\end{equation}
	Il filtro configurato in questo modo ha caratteristiche \textbf{passa-alto}.
	\subsubsection{Filtro RL}
	Risulta simile al filtro RC con la unica differenza che al posto del condensatore vi è una induttanza.
	\begin{gather}
		\bar V_L = j\omega L \bar I = j\omega L \frac{\bar E}{R+j\omega L} \\
		|\bar V_L| = E\frac{\omega L}{\sqrt{R^2 + \omega^2 L^2}} = E \frac{\omega \frac{L}{R}}{\sqrt{R+\omega ^2 \frac{L^2 }{R^2}}} \\
		\omega_L = \frac{R}{L} \implies |\bar V_L| = \frac{E}{\sqrt{2}}
	\end{gather}
	Il filtro ha caratteristiche \textbf{passa-alto} se si prende $\bar V_L$, mentre è \textbf{passa-basso} se la tensione si estrae da $R_L$.
	\begin{center}
		\begin{figure}[h]
			\centering
			\includegraphics[scale=0.55]{example-image-a}
			\caption{Circuito filtro RL}
			\label{fig:im-28}
		\end{figure}
	\end{center}
	\subsubsection{Filtro RLC Serie}
	Comprende la serie di tutti e tre i componenti, una resistenza, un condensatore e una induttanza.\\
	La corrente di maglia è data:
	\begin{equation}
		\bar I = \frac{\bar E}{R+j\omega L + \frac{1}{j\omega C}} = \frac{\bar E}{R+j(\omega L - \frac{1}{\omega C})}	
	\end{equation}
	\begin{center}
		\begin{figure}[h]
			\centering
			\includegraphics[scale=0.55]{example-image-a}
			\caption{Circuito filtro RLC}
			\label{fig:im-29}
		\end{figure}
	\end{center}
	Dati certi valori specifici di $L$ e $C$ si può raggiungere una situazione particolare in cui tensione e corrente vanno in fase, quindi la reattanza del condensatore annulla quella della induttanza. \\
	Si forma quindi un fenomeno detto \textbf{risonanza serie}.
	\begin{gather}
		\omega L = \frac{1}{\omega C} \implies \bar I = \frac{\bar E }{R} 
	\end{gather}
	Negli altri casi si ha un filtro passa-banda con una pulsazione centrale detta \textbf{pulsazione di risonanza}:
	\begin{equation}
		\omega ^\ast = \frac{1}{\sqrt{LC}}	
	\end{equation}
	\begin{center}
		\begin{figure}[h]
			\centering
			\includegraphics[scale=0.55]{example-image-a}
			\caption{Corrente in un circuito RLC}
			\label{fig:im-30}
		\end{figure}
	\end{center}
	Nel grafico in figura \ref{fig:im-30} si nota che alla pulsazione $\omega_C$ il circuito va in risonanza, e non viene effettuata alcuna attenuazione.\\\\
	La banda $B$ indica la dimensione dell'intervallo compreso fra le due pulsazioni di taglio $\omega_1$ e $\omega_2$. \\
	Il modulo della corrente si ottiene:
	\begin{equation}
		|\bar I| = \frac{E}{\sqrt{R^2+(\omega C - \frac{1}{\omega C})^2}}	
	\end{equation}
	Sapendo che $\sqrt{R^2 + (\omega L - \frac{1}{\omega C})^2} = R\sqrt{2}$ si può fare in modo di estrarre le pulsazioni di taglio:
	\begin{equation}
		\omega_{1, 2}=\sqrt{\left(\frac{R}{2L}\right)^2+\frac{1}{LC}}	\pm \frac{R}{2L}
	\end{equation}
	Da cui si trova che la prima parte della equazione rappresenta la pulsazione centrale, e quella dopo $\pm$ le pulsazioni di taglio.
	La banda risulta quindi molto semplicemente:
	\begin{equation}
		B=\frac{R}{L}	
	\end{equation}
	\subsubsection{Filtro RLC Parallelo}
	Mettendo in parallelo una induttanza e un condensatore si ottiene un filtro che elimina una determinata banda, detto "\textbf{Notch}" (o anche \textbf{Filtro Anti-risonanza} oppure ancora \textbf{Filtro Elimina-banda}).
	\begin{center}
		\begin{figure}[h]
			\centering
			\includegraphics[scale=0.55]{example-image-a}
			\caption{Circuito di un filtro Notch}
			\label{fig:im-31}
		\end{figure}
	\end{center}
	Per la frequenza $\omega$ che tende a infinito oppure zero si ha che la corrente tende a $\frac{\bar E}{R}$, mentre si ha un punto in cui essa diventa zero dove si raggiunge la frequenza $\omega^\ast$.
	\begin{center}
		\begin{figure}[h]
			\centering
			\includegraphics[scale=0.55]{example-image-a}
			\caption{Modulo della corrente per Notch}
			\label{fig:im-32}
		\end{figure}
	\end{center}
	\begin{equation}
		\bar I = \frac{\bar E}{\bar z}=\frac{\bar E}{R+\frac{j\omega L \frac{1}{j\omega C}}{j\omega L+\frac{1}{j\omega L}}}=\frac{\bar E}{R+\frac{j\omega L}{1-\omega^2LC}}	
	\end{equation}
	Quando $\omega\rightarrow 0$ la parte immaginaria si annulla e la seconda parte del denominatore diventa zero, stessa cosa per quando $\omega \rightarrow + \infty$, la frazione tende ad essere infinitesima e rimane $\frac{\bar E}{R}$. \\
	In risonanza si hanno due correnti $\bar I_C$ e $\bar I_L$ uguali e opposte che possono esser anche molto forti, quando il circuito funziona in questo modo si dice che si hanno \textbf{condizioni di antirisonanza}.\\\\
	Se si impone $1-\omega^2LC=0$ si ottiene che:
	\begin{equation}
		\omega^\ast = \frac{1}{\sqrt{LC}}	
	\end{equation}
	La banda si calcola:
	\begin{gather}
		\omega_{1, 2}=\sqrt{\frac{1}{LC}+\left(\frac{1}{2RC}\right)^2}\pm \frac{1}{2RC} \\
		B=\omega_2-\omega_1 = \frac{1}{RC}
	\end{gather}
	Un uso frequente dei filtri notch è per eliminare il rumore a 50hz presente negli apparati collegati alla corrente domestica.\\\\
	In un circuito in risonanza non si può fare in modo di utilizzare la energia dei componenti reattivi (per esempio per creare un elevatore di corrente o tensione), dato che andrebbe a "squilibrare" il fenomeno.
	\paragraph{Picchi di Risonanza o Antirisonanza Parassiti}
	\subsection{Fattore Qualità Q}
	Un altro valore importante per i filtri è il fattore di qualità, il quale misura il grado di selettività del filtro.\\\\
	Il fattore di qualità può assumere diversi significati:
	\begin{itemize}
		\item (energetico) quantità di energia immagazzinata rispetto a quella dissipata
		\item Rapporto fra centro banda e larghezza
		\item Rapporto fra tensione su uno dei componenti reattivi e totale
	\end{itemize}
	\begin{equation}
		Q:=2\pi \cdot \frac{1}{2}LI_{max}^2\cdot \frac{1}{\frac{1}{2}RI_{max}^2 T}	
	\end{equation}
	La prima parte della equazione rappresenta la energia immagazzinata da una induttanza percorsa da una corrente, mentre la seconda quella dissipata dalla resistenza in un periodo.\\\\
	Se si ha una risonanza, quindi $\omega = \frac{1}{\sqrt{LC}}=\omega^\ast$, la formula restituisce la energia di picco immagazzinata dai componenti reattivi di un filtro.
	\begin{gather}
		Q:=\frac{\omega^\ast}{B}=\frac{\omega^\ast}{\frac{R}{L}}	=\frac{\omega^\ast L}{R}\cdot \frac{I}{I}=\frac{|\bar V_L|}{|\bar E|}=V_C
	\end{gather}
	Quindi rispetto alla banda un filtro sottoposto ad una frequenza più elevata ha una qualità maggiore.
	\subsection{Funzioni di Trasferimento e di Guadagno}
	Si utilizza una funzione di trasferimento per studiare il funzionamento di filtri e circuiti. \\
	Dato un sistema con entrata $\bar X$ e uscita $\bar Y$, la funzione è data da:
	\begin{equation} \label{wow}
		\bar H(j\omega) := \frac{\bar Y(j\omega)}{\bar X(j\omega)}	
	\end{equation}
	Per trovare quindi l'uscita dato il segnale in entrata e la risposta si fa il prodotto complesso fra i due (somma di fasi, prodotto di moduli):
	\begin{equation}
		\bar Y(j\omega) = \bar X(j\omega) \cdot \bar H(j\omega)	
	\end{equation}
	Si possono distinguere diversi tipi di trasferimenti:
	\begin{itemize}
		\item $H(j\omega) = \frac{\bar I_u}{\bar I_i}$ - \textbf{Guadagno di corrente}
		\item $H(j\omega) = \frac{\bar V_u}{\bar V_i}$ - \textbf{Guadagno di tensione}
		\item $H(j\omega) = \frac{\bar V_u}{\bar I_i}$ - \textbf{Impedenza di trasferimento}
		\item $H(j\omega) = \frac{\bar I_u}{\bar V_i}$ - \textbf{Ammettenza di trasferimento}
	\end{itemize}
	Riprendendo la formula \ref{wow} si può fare in modo di scrivere numeratore e denominatore come polinomi in funzione di $j\omega$ in forma estesa ($n$ indica il numeratore, $d$ il denominatore):
	\begin{equation}
		\bar H(j\omega) = \frac{n_0 (j\omega)^n + n_1(j\omega)^{n-1} + \dots + n_{n-1}(j\omega) + n_n}{d_0(j\omega)^d+d_1(j\omega)^{d-1}+\dots+d_{d-1}(j\omega) +d_d}
	\end{equation}
	Gli ultimi termini, $n_n$ e $d_d$ sono i termini noti, e $n, d$ indicano i gradi dei polinomi. \\
	Si avranno quindi (teorema fondamentale dell'algebra) $n, d$ soluzioni denominate $z\rightarrow$ zeri (numeratore) e $p\rightarrow$ poli (denominatore). \\
	Il tutto si può scrivere quindi direttamente con le soluzioni:
	\begin{equation}
		\bar H(j\omega) = \frac{n_0}{d_0} \cdot \frac{(j\omega - z_1 ')(j\omega - z_2') \dots (j\omega - z_n')}{(j\omega - p_1')(j\omega - p_2')\dots(j\omega - p_n')}	
	\end{equation}
	Dove $\frac{n_0}{d_0}$ rappresenta il coefficiente di ordine massimo.
	\begin{definition}
		Per \textbf{zeri} delle funzioni di trasferimento indichiamo le soluzioni di $\bar n (j \omega) = 0$ cambiati di segno e gli indichiamo con i simboli: $Z_1, Z_2, \dots , Z_n$
	\end{definition}
	\begin{definition}
		Per \textbf{poli} delle funzioni di trasferimento indichiamo le soluzioni di $\bar d (j \omega) = 0$ cambiati di segno e gli indichiamo con i sombli: $P_1, P_2, \dots , P_n$
	\end{definition}
	\begin{definition}
		Definiamo \textbf{costanti di trasferimento} $k_0 = \frac{n_0}{d_0}$
	\end{definition}
	Tornando all'equazione, invertendo i segni si ha:
	\begin{equation}
		\bar H(j\omega) = \frac{n_0}{d_0} \cdot \frac{(j\omega + z_1)(j\omega + z_2) \dots (j\omega + z_n)}{(j\omega + p_1)(j\omega + p_2)\dots(j\omega + p_n)}	
	\end{equation}
	Si possono avere casi in cui uno o più poli o zeri si trovino nell'origine, quindi $p=0$ oppure $z=0$. \\
	Riscrivendo la stessa forma con i termini al denominatore, si fa in modo di estrarre poli e zeri nulli e di portarli in un termine esterno (altrimenti si creano divisioni impossibili).
	\begin{equation} \label{eq-stramba}
		\bar H(j\omega) = H_0 \frac{\left(\frac{j\omega}{z_1}+1\right)\left(\frac{j\omega}{z_2}+1\right)\dots\left(\frac{j\omega}{z_n}+1\right)}{\left(\frac{j\omega}{p_1}+1\right)\left(\frac{j\omega}{p_2}+1\right)\dots\left(\frac{j\omega}{p_d}+1\right)}(j\omega )^m
	\end{equation}
	L'ultimo termine contiene tutti i poli e zeri nulli, il numero $m$ si ottiene facendo la differenza del numero di zeri e poli all'origine.\\\\
	Una qualunque funzione di trasferimento si può esprimere in questa forma, i termini possono essere reali o complessi coniugati a coppie.
	\begin{exmp}
		Si consideri un filtro RC usato come passa-basso (fig.\ref{fig:im-33}).
		\begin{center}
			\begin{figure}[h]
				\centering
				\includegraphics[scale=0.55]{example-image-a}
				\caption{Circuito dell'esempio}
				\label{fig:im-33}
			\end{figure}
		\end{center}
	\end{exmp}
	\subsubsection{Zeri delle Fuzioni di Trasferimento}
	\subsubsection{Poli delle Funzioni di Trasferimento}
	\subsubsection{Costanti di Trasferimento}
	\subsubsection{Come Disegnare una Funzione di Trasferimento}
	\subsection{Decibel e Scala Logaritmica}
	\subsubsection{Guadagno di Potenza}
	\subsubsection{Logaritmi Fondamentali}
	\subsection{Poli e Zeri}
	\subsubsection{Reali e Distinti}
	\subsubsection{Reali Coincidenti}
	\subsubsection{Complessi e Coniugati}
	\subsection{Pulsazione di Risonanza e Fattore di Smorzamento}
\end{document}